\subsection{Parameterdarstellung einer Geraden}
\label{subsec:analytic_geometry_parametric_line_equations}

Ein Aufpunkt verankert die Gerade im Koordinatensystem; ein Richtungsvektor
legt ihren Verlauf fest. Der Parameter steuert anschließend die Bewegung in
beide Richtungen entlang der Geraden.

\begin{DefinitionBox}{Parameterdarstellung einer Geraden}{analytic_geometry_parametric_line}
  Eine Gerade $g$ durch den Punkt $P$ mit Ortsvektor $\vec p$ und mit
  Richtungsvektor $\vec a\ne\vec 0$ besitzt die
  \emph{Parameterdarstellung} \EnglishTerm{parametric equation}
  \[
    g:\quad \vec x=\vec p+\lambda\vec a,
    \qquad \lambda\in\R.
  \]
  Dabei heißt $\vec p$ \emph{Stützvektor} \EnglishTerm{support vector} und
  $\vec a$ \emph{Richtungsvektor} \EnglishTerm{direction vector}.
\end{DefinitionBox}

\noindent
\begin{minipage}[c]{0.34\linewidth}
  \centering
  \begingroup
    \def\LineUnit{0.60}
    \def\LineSupportX{2}
    \def\LineSupportY{2}
    \def\LineDirectionX{1}
    \def\LineDirectionY{-1}
    \def\LineLambdaMin{-2}
    \def\LineLambdaMax{2}
    \CachedTikzFigure{content/resources/chapters/analytic_geometry/tikz/parameterized_line_representation.tex}
  \endgroup
\end{minipage}\hfill
\begin{minipage}[c]{0.62\linewidth}
  \begin{KeyPointBox}{Bedeutung des Parameters}{analytic_geometry_line_parameter}
    Jeder Wert von $\lambda$ liefert einen Punkt der Geraden. Umgekehrt
    besitzt jeder Punkt auf $g$ einen passenden Parameterwert:
    \[
      \begin{array}{c@{\quad}l}
        \lambda=0 &: \vec x=\vec p,\\
        \lambda>0 &: \text{Richtung }\vec a,\\
        \lambda<0 &: \text{Gegenrichtung}.
      \end{array}
    \]
  \end{KeyPointBox}
\end{minipage}

\begin{ExampleBox}{Punkte einer Geraden bestimmen}{analytic_geometry_line_points_example}
  Für
  \[
    g:\quad
    \vec x=\begin{pmatrix}2\\2\end{pmatrix}
    +\lambda\begin{pmatrix}1\\-1\end{pmatrix}
  \]
  bestimmen wir drei Punkte durch gezielte Wahl des Parameters:
  \[
    \begin{array}{c@{\qquad}c@{\qquad}c}
      \lambda=-1 & \lambda=0 & \lambda=1\\[0.4em]
      \Cloze{3.5cm}{\vec x=\begin{pmatrix}1\\3\end{pmatrix}}
      &\Cloze{3.5cm}{\vec x=\begin{pmatrix}2\\2\end{pmatrix}}
      &\Cloze{3.5cm}{\vec x=\begin{pmatrix}3\\1\end{pmatrix}}
    \end{array}
  \]
  Die drei Punkte liegen auf derselben Geraden; $\lambda$ wirkt als
  Schieberegler entlang des Richtungsvektors.
\end{ExampleBox}
